\documentclass[10pt,a4paper]{amsart}
\usepackage[margin=19mm]{geometry}
\usepackage{amsmath,amssymb,mathtools}
\usepackage[scr=rsfs,cal=euler]{mathalfa}
\usepackage{xfrac}
\usepackage[unicode,hidelinks,bookmarks=false]{hyperref}
\newcommand{\ie}{i.e.\ }
\newcommand{\cf}{cf.\ }
\newcommand{\resp}{respectively\ }
\newcommand{\Subsection}{\subsection}
\providecommand{\nobreakdash}{\nobreak\hbox{-}}
\setlength{\emergencystretch}{2em}
\multlinegap0pt
\usepackage{amssymb}
\usepackage{xcolor}
\usepackage{mathtools}
\usepackage{stackengine}
\usepackage{tikz}
\newtheorem{lem}{Lemma}
\newcommand{\CAiry}{{\mathcal C_3}}%\newcommand{\CAiry}{{\mathcal C_{\mathrm A}}}
\newcommand{\CbiSone}{{\M{\mathcal C}41}}
\newcommand{\CbiStwo}{{\M{\mathcal C}42}}
\newcommand{\D}{\ensuremath{\,\mathrm{d}}}
\providecommand{\Lebesgue}{\mathrm{L}} % Upright L for Lebesgue spaces
\newcommand{\M}[3]{#1_{#2\Mspacer#3}} %Print a symbol with two subscripts eg a matrix entry
\newcommand{\Mspacer}{\;} %Spacer for below Matrix display functions
\newcommand{\RR}{\mathbb R}              % The set of reals
\newcommand{\abs}[1]{\left\lvert#1\right\rvert}
\providecommand{\argdot}{{}\cdot{}}
\providecommand{\bigoh}[1]{\mathcal{O}\left(#1\right)}
\renewcommand{\geq}{\geqslant}
\newcommand{\la}{\ensuremath{\lambda}}
\renewcommand{\leq}{\leqslant}
\newcommand{\nn}{\nonumber}
\newcommand{\normp}[2]{\left\lVert#1\right\rVert_{#2}}
\newcommand{\re}{\ensuremath{\mathrm{e}}}
\newcommand{\ri}{\ensuremath{\mathrm{i}}}
\newcommand{\what}{\widehat}
\title{Interface problems of mixed spatial order}
\author{Dionyssios Mantzavinos\textsuperscript{*}, Ravindra Pethiyagoda\textsuperscript{\S} and Dave Smith\textsuperscript{\S\textdagger}}
\date{Source: arXiv:2606.05593v1 (CC BY 4.0)}
\begin{document}
\maketitle

\noindent\emph{Source and licence.} Interface problems of mixed spatial order. Version arXiv:2606.05593v1 (CC BY 4.0), distributed by arXiv under the Creative Commons Attribution 4.0 International licence, http://creativecommons.org/licenses/by/4.0/, as recorded in the arXiv licence metadata for this exact version and shown on the abstract page https://arxiv.org/abs/2606.05593v1. Full licence notice: \texttt{licenses/arxiv-2606.05593-CC-BY-4.0.txt}.\par\smallskip

\noindent\textit{Editorial scope.} Counted unit: the article's lem lem:biSAiry.extraTermsVanish (source0.tex lines 1873--1875). The article's complete original proof of it is delivered with it (source0.tex lines 1877--1990, one \texttt{proof} environment of 114 lines). No mathematics is written, repaired or re-derived: the statement and the proof are the article's own lines, in the article's own notation, and no retained line is substituted or re-flowed.  Retained uncounted as the source-local context the proof needs: lines 806--808; lines 1685--1729; lines 1731--1858; lines 1862--1870.  Nothing was omitted from any retained span, so the package carries no omission record.  No retained line contains a citation, so the wrapper carries no bibliography and no bibliography entry.  No retained line contains a figure environment, an image-inclusion command or drawing code, so no image is delivered and the package declares no asset dependency.  Editorial formatting only, in addition: an AMS \texttt{amsart} preamble that restates the article's own declarations and macros used by the retained lines, namely the environments \texttt{lem} on separate counters, together with the standard packages those lines need.  The retained environments print in this document as Lemma 1 in order of appearance, whereas the article prints the same items with the numbering of its own full document.  The complete native source of the article as distributed in the arXiv e-print for this exact version is bundled unchanged as \texttt{sources/202128/source0.tex}.
\begin{equation} \label{eqn:lSAiry.defnCRminus}
  \mathcal C_R^- := \left\{R\re^{\ri\theta}: -\pi < \theta < 0\right\}.
\end{equation}

\begin{align}
  2\pi v(x, t)
  &=
  \int_{\RR} \re^{\ri kx+\ri k^3 t} \, \what v_0(k) \D k
  +
  \int_{\CAiry} \re^{\ri kx+\ri k^3 t} \left[\alpha \, \what v_0(\alpha k) + \alpha^2 \, \what v_0(\alpha^2 k)\right] \D k
  \nn\\
  &\qquad
  +
  \left(\alpha^2-1\right)
  \int_{\widetilde{\RR}} \re^{\ri\alpha \rho(\lambda) x+\ri\lambda t}
  \frac{\rho(\lambda)}{\Delta(\la)}
  \left[\what u_0(-\ri\gamma(\lambda)) - \what u_0(\gamma(\lambda))\right] \D\lambda
  \nn\\
  &\qquad
  +
  \int_{\widetilde{\RR}} \re^{\ri\alpha \rho(\lambda) x+\ri\lambda t}
  \frac{\gamma(\lambda)}{\Delta(\la)}
  \Big\{\left[\alpha^2 \left(1+\ri\right) \rho(\lambda) - 2\gamma(\lambda)\right] \what v_0(\rho(\lambda))
  \nn\\*
  &\hspace{17em}
  -
  \left[\left(1+\ri\right) \rho(\lambda) - 2\gamma(\lambda)\right]
  \what v_0(\alpha^2 \rho(\lambda))
  \Big\} \D\lambda
  \nn\\
  &\qquad
  -
  \left(\alpha^2-1\right)
  \int_{\widetilde{\RR}} \re^{\ri\alpha \rho(\lambda) x}
  \frac{\rho(\lambda)}{\Delta(\la)} \left[\what u(-\ri\gamma(\lambda), t) - \what u(\gamma(\lambda), t)\right] \D\lambda
  \nn\\
  &\qquad
  -
  \int_{\widetilde{\RR}} \re^{\ri\alpha \rho(\lambda) x}
  \frac{\gamma(\lambda)}{\Delta(\la)}
  \Big\{\left[\alpha^2 \left(1+\ri\right) \rho(\lambda) - 2\gamma(\lambda)\right] \what v(\rho(\lambda), t)
  \nn\\*
  &\hspace{17em}
  -
  \left[\left(1+\ri\right) \rho(\lambda) - 2\gamma(\lambda)\right]
  \what v(\alpha^2 \rho(\lambda), t)
  \Big\} \D\lambda.
  \label{eqn:biSAiry.v.withExtraTerms}
\end{align}

\begin{align}
  2\pi u(x, t)
  &=
  \int_{\RR} \re^{\ri kx+\ri k^4t} \, \what u_0(k) \D k
  -
  \int_{\CbiSone} \re^{\ri kx+\ri k^4t} \left[\left(1+i\right) \what u_0(-\ri k) - \ri \, \what u_0(-k) \right] \D k
  \nn\\
  &\quad
  -
  \int_{\CbiStwo} \re^{\ri kx+\ri k^4t} \left[\left(1-\ri\right) \what u_0(\ri k) + \ri \, \what u_0(-k) \right] \D k
  \nn\\
  &\quad
  -
  \frac{1-\ri}{2}
  \int_{\widetilde{\RR}} \re^{-\gamma(\lambda) x+\ri\lambda t}
  \left(\alpha-1\right) \rho(\lambda)
  \left[\frac{\rho(\la)}{\gamma(\la)}-\alpha^2\right]
  \left[\what u_0(\gamma(\lambda)) - \what u_0(-\ri\gamma(\lambda))\right]
  \frac{\D\lambda}{\Delta(\la)}
  \nn\\
  &\quad
  -
  \frac{1-\ri}{2}
  \int_{\widetilde{\RR}} \re^{-\gamma(\lambda) x+\ri\lambda t}
  \bigg\{
  \left[
  \left(1+\ri\right) \rho(\lambda) \left(\gamma(\lambda) + \rho(\lambda)\right) - \left(3-\ri\right) \gamma(\lambda)^2
  \right]
  \what v_0(\alpha^2 \rho(\lambda))
  \nn\\*
  &\hspace*{4cm}
  -
  \left[
  \alpha \left(1+\ri\right) \rho(\lambda) \left(\alpha \gamma(\lambda)+\rho(\lambda)\right) -\left(3-\ri\right)\gamma(\lambda)^2
  \right]
  \what v_0(\rho(\lambda))
  \bigg\}
  \frac{\D\lambda}{\Delta(\la)}
  \nn\\
  &\quad
  -
  \frac{1+\ri}{2}
  \int_{\widetilde{\RR}} \re^{-\ri\gamma(\lambda) x+\ri\lambda t}
  \left(\alpha-1\right) \rho(\lambda)
  \left[\frac{\ri\rho(\la)}{\gamma(\la)}-\alpha^2\right]
  \left[\what u_0(\gamma(\lambda)) - \what u_0(-\ri\gamma(\lambda))\right]
  \frac{\D\lambda}{\Delta(\la)}
  \nn\\
  &\quad
  -
  \frac{1+\ri}{2}
  \int_{\widetilde{\RR}} \re^{-\ri\gamma(\lambda) x+\ri\lambda t}
  \bigg\{
  \left[
  \alpha \left(1-\ri\right) \rho(\lambda) \left(\rho(\lambda)-\ri\alpha \gamma(\lambda)\right) + \left(3+\ri\right)\gamma(\lambda)^2
  \right]
  \what v_0(\rho(\lambda))
  \nn\\*
  &\hspace*{4cm}
  -
  \left[
  \left(1-\ri\right)\rho(\lambda) \left(\rho(\lambda)-\ri\gamma(\lambda)\right)
  +
  \left(3+\ri\right) \gamma(\lambda)^2
  \right]
  \what v_0(\alpha^2 \rho(\lambda))
  \bigg\}
  \frac{\D\lambda}{\Delta(\la)}
  \nn\\
  &\quad
  +
  \frac{1-\ri}{2}
  \int_{\widetilde{\RR}} \re^{-\gamma(\lambda) x}
  \left(\alpha-1\right) \rho(\lambda)
  \left[\frac{\rho(\la)}{\gamma(\la)}-\alpha^2\right]
  \left[ \what u(\gamma(\lambda), t) - \what u(-\ri\gamma(\lambda), t)\right]
  \frac{\D\lambda}{\Delta(\la)}
  \nn\\
  &\quad
  +
  \frac{1-\ri}{2}
  \int_{\widetilde{\RR}} \re^{-\gamma(\lambda) x}
  \bigg\{
  \left[
  \left(1+\ri\right) \rho(\lambda) \left(\gamma(\lambda) + \rho(\lambda)\right) - \left(3-\ri\right) \gamma(\lambda)^2
  \right]
  \what v(\alpha^2 \rho(\lambda), t)
  \nn\\*
  &\hspace*{3.5cm}
  -
  \left[
  \alpha \left(1+\ri\right) \rho(\lambda) \left(\alpha \gamma(\lambda)+\rho(\lambda)\right) -\left(3-\ri\right)\gamma(\lambda)^2
  \right]
  \what v(\rho(\lambda), t)
  \bigg\}
  \frac{\D\lambda}{\Delta(\la)}
  \nn\\
  &\quad
  +
  \frac{1+\ri}{2}
  \int_{\widetilde{\RR}} \re^{-\ri\gamma(\lambda) x}
  \left(\alpha-1\right) \rho(\lambda)
  \left[\frac{\ri\rho(\la)}{\gamma(\la)}-\alpha^2\right]
  \left[ \what u(\gamma(\lambda), t) - \what u(-\ri\gamma(\lambda), t)\right]
  \frac{\D\lambda}{\Delta(\la)}
  \nn\\
  &\quad
  +
  \frac{1+\ri}{2}
  \int_{\widetilde{\RR}} \re^{-\ri\gamma(\lambda) x}
  \bigg\{
  \left[
  \alpha \left(1-\ri\right) \rho(\lambda) \left(\rho(\lambda)-\ri\alpha \gamma(\lambda)\right) + \left(3+\ri\right)\gamma(\lambda)^2
  \right]
  \what v(\rho(\lambda), t)
  \nn\\*
  &\hspace*{3.5cm}
  -
  \left[
  \left(1-\ri\right)\rho(\lambda) \left(\rho(\lambda)-\ri\gamma(\lambda)\right)
  +
  \left(3+\ri\right) \gamma(\lambda)^2
  \right]
  \what v(\alpha^2 \rho(\lambda), t)
  \bigg\}
  \frac{\D\lambda}{\Delta(\la)}.
  \label{eqn:biSAiry.u.withExtraTerms}
\end{align}

\begin{equation} \label{eqn:biSAiry.im-ineq}
  \Im(\gamma(\lambda)) \leq 0,
  \quad
  \Im(\ri\gamma(\lambda)) \geq \frac1{\sqrt2}\abs\la \geq0,
  \quad
  \Im(-\gamma(\lambda)) \geq 0,
  \quad
  \Im(-\ri\gamma(\lambda)) \leq \frac{-1}{\sqrt2}\abs\la \leq0,
\end{equation}

\begin{lem} \label{lem:biSAiry.extraTermsVanish}
  All terms on the right of equations~\eqref{eqn:biSAiry.v.withExtraTerms} and~\eqref{eqn:biSAiry.u.withExtraTerms} which contain $\widehat u(\argdot,t)$ or $\widehat v(\argdot,t)$ evaluate to zero.
\end{lem}

\begin{proof}
  Note first that, provided $\abs\la>2^{12}$
  \begin{align*}
    \abs{\Delta(\la)} &= \sqrt6 \abs{\gamma(\la)\rho(\la)} \abs{ -\ri\alpha^2\gamma(\la)^2 + (1-\ri)\gamma(\la)\rho(\la) + \alpha\rho(\la)^2 } \\
    &\geq \sqrt6 \abs\la^{\frac14+\frac13}\left( \abs\la^{\frac23} - \sqrt2\abs\la^{\frac14+\frac13} - \abs\la^{\frac12} \right) \\
    \geq \abs\la^{\frac54}.
  \end{align*}

  Integrating by parts, we find that
  \[
    \widehat u(-\ri\gamma(\la),t) = \frac{u(0,t)}{\gamma(\la)} + \frac1{\gamma(\la)} \int_0^\infty \re^{-\gamma(\la)y}u_x(y,t)\D y,
  \]
  so
  \begin{align*}
    \abs{\widehat u(-\ri\gamma(\la),t)}
    &\leq \frac1{\abs{\gamma(\la)}} \left( \abs{u(0,t)} + \int_0^\infty \re^{y\Im(-\ri\gamma(\la))} \abs{u_y(y,t)} \D y \right) \\
    &\leq \frac1{\abs{\gamma(\la)}} \left( \normp{u(t)}{\Lebesgue^\infty_x(0,\infty)} + \normp{u_x(t)}{\Lebesgue^1_x(0,\infty)} \right),
  \end{align*}
  where we have used (the weaker version of) the latter of inequalities~\eqref{eqn:biSAiry.im-ineq} to justify the final inequality.
  Similarly,
  \begin{align*}
    \abs{\widehat u(\gamma(\la),t)}
    &\leq
    \frac1{\abs{\gamma(\la)}} \left( \normp{u(t)}{\Lebesgue^\infty_x(0,\infty)} + \normp{u_x(t)}{\Lebesgue^1_x(0,\infty)} \right),
    \\
    \abs{\widehat v(\rho(\la),t)}
    &\leq
    \frac1{\abs{\rho(\la)}} \left( \normp{v(t)}{\Lebesgue^\infty_x(0,\infty)} + \normp{v_x(t)}{\Lebesgue^1_x(0,\infty)} \right),
    \\
    \abs{\widehat v(\alpha^2\rho(\la),t)}
    &\leq
    \frac1{\abs{\rho(\la)}} \left( \normp{v(t)}{\Lebesgue^\infty_x(0,\infty)} + \normp{v_x(t)}{\Lebesgue^1_x(0,\infty)} \right).
  \end{align*}
  The above inequalities may be summarised as, uniformly in $\theta\in[-\pi,0]$,
  \begin{align*}
    \widehat u\left(\gamma\left(R\re^{\ri\theta}\right),t\right),\; \widehat u\left(-\ri\gamma\left(R\re^{\ri\theta}\right),t\right) &= \bigoh{R^{\frac{-1}4}},
    & &\mbox{as } R\to\infty \\
    \widehat v\left(\rho\left(R\re^{\ri\theta}\right),t\right),\; \widehat v\left(\alpha^2\rho\left(R\re^{\ri\theta}\right),t\right) &= \bigoh{R^{\frac{-1}3}},
    & &\mbox{as } R\to\infty.
  \end{align*}

  For $R$ greater than the deformation in $\widetilde\RR$ away from zero, let $\mathcal C_R^-$ be the contour defined by equation~\eqref{eqn:lSAiry.defnCRminus}.
  By Cauchy's theorem applied to the region bounded between $\widetilde\RR$ and $\mathcal C_R^-$,
  \[
    \int_{\widetilde\RR} \re^{\ri\alpha\rho(\la)x} \frac{\rho(\la)}{\Delta(\la)} \widehat u(-\ri\gamma(\la),t) \D\la = \lim_{R\to\infty}\int_{\mathcal C_R^-} \re^{\ri\alpha\rho(\la)x} \frac{\rho(\la)}{\Delta(\la)} \widehat u(-\ri\gamma(\la),t) \D\la,
  \]
  and equivalent contour deformation identities hold for all the other terms of interest.
  Therefore, it is sufficient to show that the limit on the right is zero and equivalent limits for the other terms are each zero.

  For $R>0$ sufficiently large, we use the above inequalities to bound
  \begin{align}
    \notag
    \abs{ \int_{\mathcal C_R^-} \re^{\ri\alpha\rho(\la)x} \frac{\rho(\la)}{\Delta(\la)} \widehat u(-\ri\gamma(\la),t) \D\la }
    &\leq
    \abs{ \int_{-\pi}^0 \re^{\ri\alpha\rho\left(R\re^{\ri\theta}\right)x} \frac{\rho\left(R\re^{\ri\theta}\right)}{\Delta\left(R\re^{\ri\theta}\right)} \widehat u\left(-\ri\gamma\left(R\re^{\ri\theta}\right),t\right) R \D \theta } \\
    \label{eqn:biSAiry.extraTermsVanish.Proof1}
    &\leq
    \int_{-\pi}^0 \abs{\re^{\ri\alpha\rho\left(R\re^{\ri\theta}\right)x}} \abs{\frac{\rho\left(R\re^{\ri\theta}\right)}{\Delta\left(R\re^{\ri\theta}\right)}} \abs{\widehat u\left(-\ri\gamma\left(R\re^{\ri\theta}\right),t\right)} R \D \theta \\
    \notag
    &\leq
    \frac{R^{\frac13}}{R^{\frac54}}\bigoh{R^{\frac{-1}4}} R \int_{-\pi}^0 \re^{-x\Im\left( \alpha\rho\left( R\re^{\ri\theta} \right) \right)} \D\theta \\
    \notag
    &\leq
    \bigoh{R^{\frac{-1}6}} \int_{-\pi}^0 \re^{-xR^{\frac13} \Im\left( \alpha\rho\left( \re^{\ri\theta} \right) \right)} \D\theta.
  \end{align}
  Now, for $\phi = (\pi-\theta)/3$,
  \(
    \Im\left( \alpha\rho\left( \re^{\ri\theta} \right) \right)
    % = \sin\left( \frac{\pi-\theta}3 \right)
    = \sin(\phi),
  \)
  so
  \[
    \int_{-\pi}^0 \re^{-xR^{\frac13} \Im\left( \alpha\rho\left( \re^{\ri\theta} \right) \right)} \D\theta
    =
    3 \int_{\frac\pi3}^{\frac{2\pi}3} \re^{-xR^{\frac13} \sin(\phi)} \D\phi.
  \]
  To obtain a bound on this integral, we use the convexity inequality
  \[
    \frac\pi3 \leq \phi \leq \frac{2\pi}3
    \Rightarrow
    \frac{3\sqrt3}{4\pi} \phi \leq \sin(\phi)
  \]
  to obtain
  \[
    3 \int_{\frac\pi3}^{\frac{2\pi}3} \re^{-xR^{\frac13} \sin(\phi)} \D\phi
    \leq
    3 \int_{\frac\pi3}^{\frac{2\pi}3} \re^{-xR^{\frac13} \frac{3\sqrt3}{4\pi} \phi} \D\phi
    \leq
    R^{\frac{-1}3} \frac{4\pi}{x\sqrt3} \left[ \re^{-xR^{\frac13} \frac{\sqrt3}{2}} - \re^{-xR^{\frac13} \frac{\sqrt3}{4} \phi} \right].
  \]
  Hence, overall,
  \[
    \abs{ \int_{\mathcal C_R^-} \re^{\ri\alpha\rho(\la)x} \frac{\rho(\la)}{\Delta(\la)} \widehat u(-\ri\gamma(\la),t) \D\la }
    = \bigoh{R^{\frac{-1}2}},
  \]
  so
  \[
    \int_{\widetilde\RR} \re^{\ri\alpha\rho(\la)x} \frac{\rho(\la)}{\Delta(\la)} \widehat u(-\ri\gamma(\la),t) \D\la = 0.
  \]

  We could have used the stronger convexity bound $\frac{\sqrt3}2\leq\sin(\phi)$ when evaluating the final integral and obtained faster decay, but analogous bounds are not available for all of the integrals from equation~\eqref{eqn:biSAiry.v.withExtraTerms}, so we have given the weaker but more broadly adaptable version of the argument.
  The above argument can be summarised as the application of the asymptotic bounds
  \begin{align*}
    \int_{-\pi}^0 \abs{\re^{\ri\alpha\rho\left(R\re^{\ri\theta}\right)x}} \D\theta &= \bigoh{R^{\frac{-1}3}}, \\
    \abs{\frac{\rho\left(R\re^{\ri\theta}\right)}{\Delta\left(R\re^{\ri\theta}\right)}} &= \bigoh{R^{\frac13}}\bigoh{R^{\frac{-5}4}},
    & \mbox{ uniformly in } \theta\in[-\pi,0], \\
    \abs{\widehat u\left(-\ri\gamma\left(R\re^{\ri\theta}\right),t\right)} &= \bigoh{R^{\frac{-1}4}},
    & \mbox{ uniformly in } \theta\in[-\pi,0], \\
    R &= \bigoh{R},
  \end{align*}
  to the integral~\eqref{eqn:biSAiry.extraTermsVanish.Proof1}.
  It is straightforward to analyse the other integrals via analogous bounds on their corresponding components.
\end{proof}

\end{document}
